\section{Discussion}
\label{sec:discussion}

\subsection{Bounded verification evidence}
\label{subsec:discussion-bounded-verification}

RQ1 provides bounded evidence that no declared property was violated in the
completed valid-model configurations. The property suite covers distinct
obligations, including replay protection, destination authorization, decision
consistency, value preservation, quorum enforcement, terminal-state
irreversibility, and release after timeout.

The evidential scope is determined by the explored bounds. Alloy completed all
three predefined profiles, whereas TLC-large was fully censored by the frozen
timeout policy. Those timeouts are not property violations; they delimit the
configurations for which bounded verification completed under the available
resource budget.

Accordingly, H1 supports only the absence of observed violations in completed
bounded configurations. It does not extend to unexplored states, censored
TLC-large executions, or an unbounded correctness claim.

TLC and Alloy provide complementary corroboration rather than interchangeable
proofs. Their different formal semantics and state representations prevent
agreement across completed profiles from implying identical explored spaces or
equivalent proofs.

\subsection{Mutation sensitivity and property adequacy}
\label{subsec:discussion-mutation}

RQ2 complements RQ1 by asking whether the declared properties react when
selected protocol controls are removed. All ten predefined scientific mutants
were detected by their designated target properties, producing a mutation score
of 1.0. This makes the valid-model results more informative because controlled
defects in replay protection, receipt-dependent credit, terminal decisions,
timeout release, and quorum enforcement trigger the expected properties.

Tool-specific encodings can express the same defect intention through different
properties. For example, commit after abort is associated with
\texttt{TerminalState\allowbreak Irreversibility} in Alloy and
\texttt{Decision\allowbreak Consistency} in TLA+. The mutant-level conclusion
therefore concerns targeted defect detection, not identical property
formulations across tools.

The perfect mutation score remains restricted to the predefined catalogue. It
does not establish completeness against arbitrary implementation defects,
modelling errors, or other cross-shard failure modes.

RQ1 and RQ2 thus support complementary claims: the declared properties hold in
the completed valid-model configurations, and the evaluated property suite is
sensitive to the selected protocol controls when they are deliberately
weakened.

\subsection{Implementation-model trace conformance}
\label{subsec:discussion-conformance}

RQ3 adds implementation-facing evidence by replaying observable Java traces
against the TLA+ abstraction. The valid catalogue was accepted and the
deliberately corrupted catalogue rejected as expected, connecting the formal
model with the implementation behaviors exercised by the frozen scenarios.

The negative controls strengthen this result because rejection was checked at
the expected diagnostic point rather than treated as a generic replay failure.
This provides evidence that deliberately introduced mismatches are localized at
the intended abstract-action boundary.

The result remains limited to the selected observable abstraction, scenario
catalogue, negative mutations, and frozen seeds. It therefore supports bounded
implementation-model trace conformance rather than a statement about every
possible Java execution.

In particular, this evidence must not be interpreted as a formal
refinement proof or as behavioral equivalence between the Java
implementation and the TLA+ specification. Establishing refinement
would require a stronger semantic relation covering the full relevant
transition systems, whereas the present experiment validates
correspondence for the predefined observable traces.

RQ3 therefore complements internal model checking and mutation analysis by
confronting the formal abstraction with observable implementation executions.

\subsection{Verification cost and scalability}
\label{subsec:discussion-cost}

RQ4 shows that bounded verification is constrained by computational cost as
well as logical outcome. Measurements increase within both formal tools, but
H4 remains unconfirmed in its preregistered nonlinear form.

For Alloy, all three ordered profiles completed, and median elapsed time and
memory increased monotonically from small to medium and large. The Spearman
association of 1.0 for both metrics is consistent with that ordering, but three
ordinal points cannot identify a functional growth law.

TLC exposes a sharper feasibility boundary. The small and medium profiles
completed with increases in elapsed time, memory, and explored state-space size,
whereas all seventy TLC-large measured runs reached the frozen 1800-second
timeout. Complete censoring precludes a completed-run median for the large
profile and a three-level characterization comparable to Alloy. The timeout is
therefore a valid censored observation, not an experimental failure, and
retaining it makes the evaluated verification boundary explicit.

Within the TLC medium-bound configurations, enabled fault profiles also changed
median time and memory across normal, insufficient-quorum, replay, and
timeout-oriented cases, indicating that verification cost depends on behavioral
structure as well as nominal bound size. These differences do not identify a
universal cost ordering, but they show that enabled behavior is part of the
verification envelope.

The evidence supports only a descriptive conclusion that bounded verification
cost increases for the evaluated configurations. It establishes neither a
nonlinear nor an asymptotic growth law. TLC and Alloy also use different
semantics, exploration mechanisms, and internal representations, so absolute
time, memory, or state-space measurements cannot support intrinsic
tool-superiority claims.

\subsection{Cross-layer evidence triangulation}
\label{subsec:discussion-triangulation}

The four research questions are intended to provide complementary evidence
rather than four independent demonstrations of protocol correctness. Each
addresses a different weakness in the validation argument.

RQ1 evaluates whether the declared properties survive the bounded
interleavings explored for valid formal models. Because absence of a
counterexample does not demonstrate sensitivity to relevant defects, RQ2
deliberately weakens selected protocol controls and checks whether their
designated properties detect the resulting violations.

Formal verification and mutation detection still concern formal
representations. RQ3 adds implementation-facing evidence by checking bounded
conformance of observable Java traces against the TLA+ abstraction, including
deliberately corrupted executions. This connects the formal layer with the
implementation without implying refinement or behavioral equivalence.

RQ4 makes the computational boundary explicit by characterizing verification
cost as the evaluated bounds become more demanding. The observed increase in
cost and complete TLC-large censoring show that the formal evidence is obtained
within a finite resource envelope.

Together, the layers form an evidence chain from bounded property verification
through mutation-based property validation and bounded implementation-model
trace conformance to verification-cost characterization. The chain is
complementary rather than cumulative proof: model checking remains bounded,
mutation analysis remains catalogue-dependent, trace replay remains limited to
the chosen abstraction and scenarios, and cost characterization remains tied to
the evaluated configurations. No layer removes the limitations of another.

The combined result is therefore not a general proof of cross-shard protocol
correctness. It is a reproducible multi-layer argument under one frozen
experimental protocol, connecting formal models, executable evidence, and the
resource conditions under which that evidence was obtained while preserving the
scope limits of each layer. This integration, rather than any individual
verification technique, is the methodological contribution emphasized here.
